
\subsection{Design Considerations}

When designing an \gls{llm} assistant using small \glspl{llm}, several problems
became apparent. These include:

\subsubsection{Limited Context Window}

Small \glspl{llm} are constrained in processing large amounts of text at
once, which poses challenges when using Retrieval-Augmented
Generation (RAG). This limitation inherently restricts the model's
ability to handle complex queries or tasks requiring extensive
context, necessitating innovative strategies for context management.

\subsubsection{RAG --- Prompt Size Overhead}

\subsubsection{Tokenizer Limitations}

Small \glspl{llm} often struggle to handle complex, non-human structures such as bash
commands and flags (for example, confusing -m with -M) commonly found in CLI
commands and man pages. This limits their applicability when dealing with such
structured data, leading to challenges in understanding and generating relevant
outputs.

\subsubsection{Document Relevance Filtering}

\subsubsection{Iterative Prompt Processing}

To further reduce prompt sizes, the RAG system was customized to
process each chosen document in smaller parts, aggregating results
afterward. This approach avoids overwhelming the model with excessive
input at once and enhances both precision and processing efficiency.

\subsection{System Architecture}

Based on previously reviewed solutions, as well as small local \gls{llm} limitations,
the best approach to system architecture is to divide the whole system into
smaller components, each one responsible for one simple task. These defined
small tasks will (1) reduce the possibility of \gls{llm} halucination and (2)
simplify the testing by alloing us to test and evaluate each component in
separate. The modular design also makes future maintenance and expansion more
manageable.

% repetitive prompting >> onebigprompt due to context + halucinations

\begin{enumerate}
  \item \textbf{Natural Language vs. Shell Command Classifier}
    --- will allow the system to distinguish between user commands and
    questions, thus simplifying the interaction for the user.  This is
    critical for seamless operation; the system should intelligently
    decide whether to execute a command directly or to engage the \gls{llm}
    for a more conversational response.
  \item \textbf{Information Retriever} --- for specific tasks
    related to information retrieval, the RAG-enabled component will be
    designed. It should be able to retrieve information from manuals or
    the web to provide accurate answers. This component will leverage
    a knowledge base (local documents, indexed websites, etc.) and
    utilize retrieval techniques to provide contextually relevant
    information to address user inquiries.
  \item \textbf{Task splitter} ---
    this component will split complex tasks into smaller ones, thus
    dividing the work for other nodes. This is inspired by multi-agent
    architecture and will follow similar principles.  It will be
    responsible for taking a user's request (after classification) and
    breaking it down into sequential or parallel sub-tasks that can be
    handled by the Information Retriever or other subsequent
    components.
\end{enumerate}

The interface will remain as close as possible to CLI with certain
TUI elements where needed.  We will prioritize a minimal,
predictable user experience.

\subsection{Development Process}

\subsubsection{Component 1. Natural Language vs. Shell Command Classifier}

\textbf{Idea.}{}
The first component is a simple proof-of-concept implementation of the proposed
idea. It uses a local \gls{llm} hosted by Ollama to distinguish between shell
commands and natural language input. If the input is a shell command, it would
execute it directly; if it's natural language, it would respond with an answer
or suggestion.\\

While using such architecture is feasible, the problem that arises is how
should the wrapper interact with the shell. It cannot attach to an existing
bash process, as it is not possible to determine when the shell is waiting for
user input (e.g.~after running a TUI program) and when it is executing commands
(this will be visible when using TUI applications inside the wrapper).\\

Another approach is to create new shell process for each command, but this will
lead to loss of context (e.g.~current working directory, environment variables,
etc.) and will make the wrapper less useful.

\textbf{Testing. }{}

\subsubsection{Component 2. Information retriever}

\textbf{Idea.}{}
This component is responsible for retrieving information from various sources
to answer user queries. It's built around a Retrieval-Augmented Generation
(RAG) architecture, meaning it combines information retrieval from a local
knowledge base with \gls{llm} generation to produce informative and accurate
responses. The task of this module is to summarize the relevant context and
pass it to user or requesting node, which then generates a response based on
the retrieved information and its own knowledge.\\

\textbf{RAG Overhead and Proposed Solutions.}{}
One of the key challenges faced while implementing RAG in small \glspl{llm}
is the substantial prompt size it generates. This overhead negatively
impacts performance, precision, and latency. To address this, the
following two-step approach has been adopted:

\begin{enumerate}
  \item \textbf{Document Relevance Filtering:} An additional step is
    implemented where the assistant selects the most relevant
    documents based on easily identifiable beginnings (e.g., command
    names and short descriptions in manual pages). This is designed
    to optimize the initial retrieval process and ensure minimal
    context is sent to the \gls{llm}.
  \item \textbf{Iterative Prompt Processing:} To further optimize,
    the RAG system processes documents in small, manageable parts
    rather than in full. The results are aggregated after processing,
    reducing latency and improving precision.
\end{enumerate}

\textbf{Testing. }{}

\textbf{Testing. }{}

\textbf{Retrieval Module.}{}
The task of this module is, provided with a path to a folder with resources
(man pages, in our case), assuming that each document in it can be assigned one
topic, to (1) search the data and return top-k document identifiers (file
paths). It will use a locally run embedding model to (2) index these documents
and enable contextual search.\\

The interface it should provide will consist of mainly three functions:

\begin{enumerate}
  \item \texttt{build (db\_path, index\_path\_within\_db,
    OpenAIAPI-Client, model)} --- will build
    an index at the specified \texttt{db\_path}, and save it, along
    with the model used for
    embedding creation (as \texttt{config.json}) into
    \texttt{db\_path/index\_path\_within\_db}.
    \texttt{index\_path\_within\_db} is a non-mandatory parameter,
    and its default value is
    \texttt{.index}. For creating embeddings, it will use
    \texttt{OpenAIAPI-Client} with the specified model.

  \item \texttt{check (db\_path, index\_path\_within\_db)} --- will
    check the availability of the index
    within \texttt{db\_path}.

  \item \texttt{search (db\_path, index\_path\_within\_db,
    OpenAIAPI-Client)} --- will conduct a search
    on \texttt{db\_path}, using the index from
    \texttt{db\_path/index\_path\_within\_db}, with
    \texttt{OpenAIAPI-Client} and the model from
    \texttt{db\_path/index\_path\_within\_db}. Calls
    \texttt{check (db\_path, index\_path\_within\_db)} --- will check
    the availability of the index
    within \texttt{db\_path} and \texttt{index\_path\_within\_db}.\\
\end{enumerate}

It should also implement its own CLI interface with \texttt{build} and
\texttt{search} commands for testing purposes.

This implementation will be using following technologies:
\begin{itemize}
  \item \textbf{Embedding model}: We will be choosing a local
    embedding model from:
    \begin{itemize}
      \item \texttt{qwen3-embedding:0.6b}
      \item \texttt{ibm/granite-embedding:125m}
    \end{itemize}
  \item \textbf{Vector database}: We will be using \texttt{FAISS} for
    indexing and searching the embeddings. We will test different
    similarity metrics (e.g.~cosine similarity, inner product) to find the
    best one for our use case.
  \item \textbf{Connection to \gls{llm}}: We will be using
    \texttt{OpenAIAPI-Client} to connect to the local
    \gls{llm} for generating responses based on the retrieved information.
\end{itemize}

\textbf{Man pages processing for embedding database.}{}

For better context management and embedding generation, the pages will be
compiled to text format (see groff -Tascii and grotty). This will allow us to
extract the most relevant information without the noise of formatting and other
non-essential elements.

\textbf{RAG implementation and information extraction by \gls{llm}.}{}

The naive approach to RAG implementation would be to retrieve the top-k
relevant documents and pass them all to the \gls{llm} in one prompt. However, this
approach generates great overhead in terms of prompt size,
negatively affecting performance, precision, and latency of small
\glspl{llm}.

To minimize RAG-induced overhead, an additional step was implemented
where the \gls{llm} selects the most relevant documents based on their
beginnings (e.g., command names and short descriptions in man pages).
This pre-filtering process is designed to improve accuracy and reduce
the computational load, ensuring that only essential information is processed.

Another optimization strategy is to process the retrieved documents in smaller,
manageable parts rather than in full. This can be implemented in two ways:
first, where the \gls{llm} processes one document at a time and aggregates results,
or second, where, before passing the documents to the \gls{llm}, they are split into
smaller parts (e.g., sections or paragraphs) and filtered by the embedding
model for relevance. The second approach is more efficient as it reduces the
number of interactions with the \gls{llm}, thus minimizing latency and improving
precision.


\subsubsection{Component 3. Task splitter}

\textbf{Idea.}{}
This component serves as the central coordinator for more complex requests. It
receives the user input and decomposes it into a series of manageable sub-tasks.
It will use a strategy inspired by multi-agent coordination principles.

\textbf{Testing. }{}

\subsection{Deployment}

To deploy the \gls{llm} node a docker image is used hosting an Ollama image. This
ensures a controlled and isolated environment for running the local \glspl{llm}
effectively.

\textbf{Service Checking Logic:} A dedicated service-checking routine
is implemented to verify that the \gls{llm} is running and available for
processing requests. This service runs at regular intervals and
includes retries and logging mechanisms to troubleshoot.

\subsection{Libraries and Technology Stack}

The development of this \gls{llm} assistant utilizes the following
libraries and technologies:

\begin{itemize}
  \item \textbf{Python Libraries:}
    \begin{itemize}
      \item \texttt{faiss-cpu}: For vector similarity search and indexing.
      \item \texttt{numpy}: Numerical operations.
      \item \texttt{openai}: Implementation of the OpenAI API client
        for \gls{llm} interactions.
    \end{itemize}
  \item \textbf{Embedding Models:}
    \begin{itemize}
      \item \texttt{qwen3-embedding}: Versions 0.6b, 1.7b, 4b.
      \item \texttt{ibm/granite-embedding}: Versions 125m, 268m.
    \end{itemize}
  \item \textbf{Vector Database:}
    \begin{itemize}
      \item \texttt{FAISS}: For fast similarity search.
    \end{itemize}
  \item \textbf{Containerization and services:}
    \begin{itemize}
      \item \texttt{Docker}: For containerized deployment.
      \item \texttt{Ollama}: For hosting local \glspl{llm}.
    \end{itemize}
\end{itemize}
